% DRAFT for ICDE 2027 R1 EAB submission.
% Author: collaboratively drafted 2026-05-28 (Claude + user).
% NOT YET MERGED into main.tex. Review and edit, then copy in.
%
% Word budget targets: abstract ~200 words; introduction ~1 page (IEEE 2-col).
% Page-budget cuts elsewhere (current 14p -> 12p limit) reserved for §4.3
% static-vs-dynamic ablation + §5 (a) Why-static-beats-dynamic (per kickoff plan).

\title{[Experiment, Analysis, and Benchmark] An Empirical Study of 3D-Pair Injection
Mechanisms in Pretrained Molecular Spectral Graph Neural Networks}

% ======================================================================
% Abstract (draft v1, ~210 words; trim to 200 in final pass)
% ======================================================================

\begin{abstract}
Pretrained graph neural networks for molecular property prediction
routinely incorporate three-dimensional structural information through
pair representations, yet the contribution of these representations to
downstream predictive performance under matched evaluation has not been
systematically audited. We conduct a controlled empirical study of three
independently motivated 3D-pair injection mechanisms---static bond-level
gating from a frozen Uni-Mol encoder (V2-T5), iterative node-to-pair
update (T6), and pair-biased multi-head attention with a per-head sigmoid
gate (T7)---within a shared ChebNet II spectral backbone trained with
NT-Xent contrastive pretraining. On seven MoleculeNet-derived benchmarks
(four classification: BBBP, BACE, ClinTox, Tox21; three regression:
FreeSolv, ESOL, Lipo) under a matched Uni-Mol scaffold-fold split protocol
with three to nine training seeds per cell, all three deep variants tie
or trail a 2D-only baseline on classification within seed variance, while
the V2-T5 + Uni-Mol pair pipeline achieves the strongest regression result
on FreeSolv. A random forest on Morgan + RDKit 2D descriptors evaluated
under a 30-seed protocol dominates BACE classification by approximately
13 AUC points, demonstrating fingerprint saturation as the ceiling on at
least one benchmark. A mechanism-level analysis of V2-T5 reveals an
aromatic-routing pattern (Pearson \( r = {+}0.66 \)) consistent across all
audited training seeds, and T7's per-head attention gate stays in the
near-identity-residual regime documented by ReZero and LayerScale across
every audited checkpoint. We additionally document and quantify a
silent unidirectional-bond featurization regression in our codebase that
drops approximately 19\% of heavy atoms from the symmetric Laplacian and
shifts apparent V2-T5 BACE performance by 7.4 AUC points; we audit
fourteen widely cited public molecular GNN repositories for the same
vulnerability. We release the audit script, all dataset processing
pipelines, baseline implementations, and per-seed checkpoints to support
reproducible cross-benchmark comparison.
\end{abstract}

% ======================================================================
% Introduction (draft v1, ~1 page IEEE 2-col estimated)
% ======================================================================

\section{Introduction}
\label{sec:intro}

Three-dimensional pair representations from pretrained molecular encoders
(notably Uni-Mol~\cite{uni-mol2023} and its descendants) have become a
standard input modality for downstream molecular property prediction.
The implicit hypothesis behind their adoption is that pair-wise geometric
information, integrated into the message-passing graph encoder, transfers
useful inductive bias from a large unlabeled chemistry corpus to small
labeled benchmarks. Concurrently, a wave of reproducibility studies has
shown that the apparent gains of deep graph encoders over classical
fingerprint-based baselines are protocol-sensitive
\cite{deng2023systematic, vanTilborg2024limitations, pocha2023atomic}.
This study asks a focused empirical question that sits between those two
literatures: when matched against strong classical and supervised-graph
controls under a single evaluation protocol, do three independently
motivated 3D-pair injection mechanisms differ predictively, and if not,
what mechanisms do they actually learn?

\paragraph{Mechanisms studied.}
We compare three injection paths within a shared ChebNet II
\cite{he2022chebnetii} spectral backbone:
(i) \emph{V2-T5}, a small MLP-projected gate (with bias-initialized
near-identity sigmoid) that multiplies each chemical-bond edge weight in
the symmetric Laplacian;
(ii) \emph{T6}, an iterative bidirectional node-to-pair update that lets
both representations refine each other at every layer; and
(iii) \emph{T7}, a pair-biased multi-head attention block with a per-head
sigmoid gate stacked on top of V2-T5. All three share the same
contrastive pretraining objective (NT-Xent across four views), the same
downstream evaluation (frozen encoder + linear probe), and the same
fingerprint MLP fusion. They differ only in how the frozen pair
representation enters the spectral filter.

\paragraph{Benchmark scope.}
Following Luo \emph{et al.} \cite{luo2024classicgnn}, we report all
variants on a single matched protocol across seven MoleculeNet-derived
benchmarks: BBBP (n=2{,}039), BACE (n=1{,}513), ClinTox (n=1{,}478,
2-task), and Tox21 (n=7{,}831, 12-task) for classification; FreeSolv
(n=642), ESOL (n=1{,}128), and Lipo (n=4{,}200) for regression. Each
training cell is averaged over three pretraining seeds with three
downstream evaluation restarts (nine measurements) where the budget
permits, and the BBBP V0/V2-T5/T6 cells additionally extend to nine
training seeds to anchor the seed-variance discussion. Classical Random
Forest with Morgan + RDKit 2D descriptors is reported under the 30-seed
protocol of Deng \emph{et al.}~\cite{deng2023systematic} alongside a
supervised D-MPNN baseline (Chemprop)~\cite{yang2019chemprop}.

\paragraph{Contributions.}
\begin{itemize}
\item \textbf{(C1) A controlled cross-benchmark audit of 3D-pair injection.}
On seven benchmarks under matched seeds and split protocol, all three
deep variants tie or trail a 2D-only baseline on every classification
endpoint within seed variance. The largest reliable separation is on
FreeSolv regression, where V2-T5 + Uni-Mol pair achieves RMSE 0.638
against Chemprop's default-hyperparameter 0.879. Random Forest dominates
BACE by approximately 13 AUC points, replicating the fingerprint-
saturation pattern of~\cite{deng2023systematic} under our setting.

\item \textbf{(C2) Mechanism-level analysis of V2-T5.} The learned bond
gate exhibits an aromatic-routing pattern: across all nine audited (BBBP,
seed) checkpoints the gate up-weights aromatic bonds and down-weights
saturated single bonds, with Pearson correlation against the aromatic-
bond indicator in $[0.6, 0.7]$. The same architecture trained on BACE
does not learn the routing, consistent with fingerprint saturation
masking any geometry-side signal.

\item \textbf{(C3) Documentation of T7's closed-gate regime.} The
per-head sigmoid gate in T7 stays within $[0.122, 0.124]$ across all nine
audited (benchmark, seed) checkpoints. This places the trained T7 in the
near-identity-residual regime documented by ReZero
\cite{bachlechner2021rezero} and LayerScale \cite{touvron2021cait}, and is
consistent with the contrastive pretraining objective providing little
direct gradient pressure on attention-output utilization.

\item \textbf{(C4) Forensic featurization audit.} An internal
unidirectional-bond featurization regression in our GraphDTA-derived
loader~\cite{nguyen2021graphdta} caused approximately 19\% of heavy atoms
on BBBP, 18\% on BACE, and 24\% on FreeSolv to receive out-degree zero in
the symmetric Laplacian. Controlling for this regression in a matched
pre-fix vs.\ post-fix BACE ablation shifts V2-T5 by 7.4 AUC points,
demonstrating that featurization choices alone can match or exceed the
gap that motivates architectural development. We audit fourteen widely
cited molecular GNN repositories for the same vulnerability and report
where it is and is not present.

\item \textbf{(C5) Reproducibility artifact.} Per the ICDE EAB requirement,
we release the full evaluation harness, all dataset preparation scripts
including the B4 featurization audit, baseline implementations, and
per-seed checkpoints. All numbers in our tables are regenerated by a
single \texttt{paper/make\_tables.py} from machine-readable per-seed JSON.
\end{itemize}

The remainder of the paper is organized as follows. Section
\ref{sec:related} reviews 3D-pair injection in molecular GNNs and
benchmark-reliability work in molecular ML. Section \ref{sec:method}
describes the shared backbone, the three injection mechanisms, and the
featurization-control protocol. Section \ref{sec:experiments} presents
the cross-benchmark audit. Section \ref{sec:discussion} discusses the
mechanism-without-metric finding, fingerprint saturation, and
limitations. Section \ref{sec:conclusion} concludes.
